\section{The Gate: What Runs Without a Board}
\label{sec:build-gate}

A project with hardware has a scheduling problem: the compiler is available constantly, the board is
available sometimes. Anything that needs the board gets deferred, and deferred checks do not happen.
The answer this project uses is a single script that runs everything that does not need hardware.

\subsection{What It Checks}

\begin{minted}{text}
$ bash scripts/gate.sh
PASS  cargo fmt --all --check
PASS  efr32mg24-hal clippy, 0 warnings
PASS  efr32mg24-pac clippy, 0 warnings
PASS  xiao-expansion-base-bsp clippy, 0 warnings
PASS  xiao-mg24-bsp clippy, 0 warnings
PASS  motion-detection clippy, 0 warnings
PASS  efr32mg24-hal: 5 examples of its own level
PASS  efr32mg24-pac: 19 examples of its own level
PASS  xiao-expansion-base-bsp: 9 examples of its own level
PASS  motion-detection: host tests (9 passed)
PASS  motion-detection: no hardware dependency
PASS  examples outside the PAC crate name no PAC type
PASS  no BSP re-exports the PAC
PASS  the HAL exposes no register block publicly
PASS  no example steals a core peripheral
VERDICT: PASS
\end{minted}

Every line is a property, not a command. That distinction matters: a script that runs
\codeinline{cargo clippy} tells you a tool ran; a line that says \emph{zero warnings in five crates}
tells you what is true.

\subsection{The Discovered-vs-Listed Rule}

The first version of this gate had a hard-coded list of crates, and the crate that most needed
checking was the one that was not on it. The version in use \emph{discovers} crates:

\begin{minted}{bash}
for dir in crates/*/ libraries/*/; do
  [ -f "${dir}Cargo.toml" ] || continue
  name=$(basename "$dir")
  ...
done
\end{minted}

A list is a snapshot of what existed when someone wrote it. Discovery cannot fall behind — which is
exactly what a check is for.

\subsection{The Rule Nobody Wrote Down}

The most valuable check in that output is the one that sounds odd: \textbf{every crate must have at
least one example of its own level}.

\begin{itemize}
    \item An API with zero examples passes every other check trivially — it compiles, it lints, and
          nothing calls it.
    \item It then rots silently: in this project a whole board module spent weeks documented and
          unbuildable, and the only symptom was that nobody used it.
    \item The check is cheap and it fails loudly, which is the entire point of a gate.
\end{itemize}

\subsection{What a Gate Cannot Do}

It cannot tell you the LED lit. This project's gate was green on a firmware that flashed successfully
and did not blink — four green steps, no light. That is not a failure of the gate; it is the boundary
of what a build-time check can know, and it is why the hardware checks in chapter~14 exist as a
separate, written procedure with its own measurements.
